\documentclass[12pt,oneside]{book}

% ============================================================
% METADATOS
% ============================================================
\newcommand{\universidad}{Universidad de Buenos Aires (UBA)}
\newcommand{\facultad}{Facultad de Derecho}
\newcommand{\materia}{Derecho Societario}
\newcommand{\titulo}{La transición del sistema de tipicidad hacia la autonomía de la voluntad}
\newcommand{\autor}{Isaac Edgar Camacho Ocampo}
\newcommand{\anio}{2026}

% ============================================================
% PAQUETES
% ============================================================
\usepackage[utf8]{inputenc}
\usepackage[T1]{fontenc}
\usepackage[spanish,es-nodecimaldot]{babel}

\usepackage[
    top=2.5cm,
    bottom=2.5cm,
    left=3cm,
    right=2.5cm,
    headheight=15pt
]{geometry}

\usepackage{amsmath}
\usepackage{amssymb}
\usepackage{mathtools}

\usepackage{graphicx}
\usepackage{float}
\usepackage{caption}
\usepackage{subcaption}

\usepackage{booktabs}
\usepackage{array}
\usepackage{longtable}
\usepackage{tabularx}

\usepackage{enumitem}

\usepackage{xcolor}
\usepackage[most]{tcolorbox}

\usepackage{fancyhdr}
\usepackage{ragged2e}

% ============================================================
% RUTAS DE IMÁGENES
% ============================================================
\graphicspath{
    {./img/}
    {./graficos/}
    {./figuras/}
}

% ============================================================
% CONFIGURACIÓN
% ============================================================
\setcounter{tocdepth}{2}

\setlength{\parindent}{0pt}
\setlength{\parskip}{0.5em}

\setlist[itemize]{
    topsep=0.3em,
    itemsep=0.2em
}

\setlist[enumerate]{
    topsep=0.3em,
    itemsep=0.2em
}

\clubpenalty=10000
\widowpenalty=10000

\pagestyle{fancy}
\fancyhf{}

\fancyhead[L]{\nouppercase{\leftmark}}
\fancyhead[R]{\materia}

\fancyfoot[C]{\thepage}

\renewcommand{\headrulewidth}{0.4pt}
\renewcommand{\footrulewidth}{0pt}

\fancypagestyle{plain}{
    \fancyhf{}
    \fancyfoot[C]{\thepage}
    \renewcommand{\headrulewidth}{0pt}
}

% ============================================================
% COMANDOS
% ============================================================
\newcommand{\abs}[1]{\left|#1\right|}
\newcommand{\norm}[1]{\left\lVert#1\right\rVert}
\newcommand{\vect}[1]{\mathbf{#1}}
\newcommand{\deriv}[2]{\frac{d#1}{d#2}}
\newcommand{\pderiv}[2]{\frac{\partial#1}{\partial#2}}

% ============================================================
% ENTORNOS / CAJAS ACADÉMICAS
% ============================================================
\newtcolorbox{definition}[1]{
    enhanced,
    breakable,
    title={Definición: #1},
    fonttitle=\bfseries,
    colback=blue!3,
    colframe=blue!50!black,
    boxrule=0.8pt,
    arc=2mm
}

\newtcolorbox{important}{
    enhanced,
    breakable,
    title={Importante},
    fonttitle=\bfseries,
    colback=yellow!8,
    colframe=orange!70!black,
    boxrule=0.8pt,
    arc=2mm
}

\newtcolorbox{example}[1]{
    enhanced,
    breakable,
    title={Ejemplo: #1},
    fonttitle=\bfseries,
    colback=green!4,
    colframe=green!40!black,
    boxrule=0.8pt,
    arc=2mm
}

\newtcolorbox{formula}[1]{
    enhanced,
    breakable,
    title={Fórmula / Regla: #1},
    fonttitle=\bfseries,
    colback=purple!4,
    colframe=purple!50!black,
    boxrule=0.8pt,
    arc=2mm
}

\newtcolorbox{exercise}[1]{
    enhanced,
    breakable,
    title={Ejercicio: #1},
    fonttitle=\bfseries,
    colback=gray!5,
    colframe=gray!60!black,
    boxrule=0.8pt,
    arc=2mm
}

\newtcolorbox{note}{
    enhanced,
    breakable,
    title={Nota},
    fonttitle=\bfseries,
    colback=white,
    colframe=gray!50,
    boxrule=0.6pt,
    arc=2mm
}

% ============================================================
% CONFIGURACIÓN FINAL: HIPERVÍNCULOS Y METADATOS PDF
% ============================================================
\usepackage[
    colorlinks=true,
    linkcolor=blue!50!black,
    citecolor=green!40!black,
    urlcolor=blue!60!black,
    pdftitle={\titulo},
    pdfauthor={\autor},
    pdfsubject={Apuntes universitarios},
    bookmarks=true
]{hyperref}

% ============================================================
% DOCUMENTO
% ============================================================
\begin{document}

% ---------------- PORTADA ----------------
\begin{titlepage}

\thispagestyle{empty}

\begin{center}

\vspace*{1cm}

{\large \textsc{\universidad}}

\vspace{0.2cm}

{\large \textsc{\facultad}}

\vspace{3cm}

{\Huge \textbf{\materia}}

\vspace{1cm}

{\LARGE \titulo}

\vspace{0.8cm}

\textit{Comprender el porqué de cada norma es el primer paso para ejercer el derecho con criterio propio.}

\vspace{1.2cm}

\rule{0.70\textwidth}{0.5pt}

\vspace{0.5cm}

{\large Apuntes de estudio}

\vspace{0.5cm}

\rule{0.70\textwidth}{0.5pt}

\vfill

{\large \autor}

\vspace{0.3cm}

{\large \anio}

\end{center}

\vfill

\begin{flushleft}
\textit{Este es un modesto aporte a los estudiantes de ``Derecho Societario'' no pretende sustituir el material oficial de la materia.}
\end{flushleft}

\end{titlepage}

% ---------------- ÍNDICE ----------------
\tableofcontents

% ============================================================
% CAPÍTULO 1
% ============================================================
\chapter[Presentación y contexto del cambio]{Presentación y contexto del cambio}

\begin{note}
La presente unidad temática recorre la metodología de la cátedra, el diagnóstico sobre la transformación del mundo del trabajo y el contexto de la reforma proyectada de la Ley General de Sociedades, como marco introductorio para el estudio del derecho societario.
\end{note}

\section{Metodología de la cátedra}

\subsection{Un curso variable según el grupo}

La docente explicó, como principio rector de su forma de enseñar, que ningún curso resulta igual al anterior, aunque se trate de la misma materia:

\begin{quote}
\textit{``Yo les cuento, yo soy medio atípica y mis cursos no son ninguno igual. Si vos viniste en el primer bimestre y venís ahora, ya la misma materia va a ser totalmente diferente. ¿Por qué? Porque ustedes son distintos.''}
\end{quote}

Según explicó, el nivel y la profundidad de cada curso se adaptan a lo que el grupo exige y está dispuesto a dar: la base del programa se mantiene, pero lo que se construye sobre esa base depende del compromiso del grupo.

\begin{important}
El aprendizaje universitario es un proceso bidireccional: el nivel de exigencia de una clase no depende únicamente del docente, sino también de la disposición del grupo a involucrarse.
\end{important}

La docente insistió en que la etapa universitaria es la última oportunidad de formarse sin costo directo, antes del ejercicio profesional:

\begin{quote}
\textit{``Ustedes van a estar en la calle, con lo cual esto es gratuito, aprovéchenlo, no se queden con la base, porque después todo lo que tengan que aprender lo van a tener que pagar.''}
\end{quote}

Remarcó que, en esta etapa final de la carrera, elegir materias ``para safar'' ya no tiene sentido: hay que elegir para formarse, porque la formación es lo que permitirá diferenciarse en un mercado profesional cada vez más competitivo.

\subsection{Forma de trabajo y uso de herramientas de inteligencia artificial}

La docente explicó que el dictado de la materia no se centrará en la exposición de conceptos que hoy están disponibles gratuitamente en cualquier buscador o modelo de lenguaje:

\begin{quote}
\textit{``Si yo voy a dar un concepto hoy, honestamente no me necesitan, porque hoy entran al chat GPT y les va a dar lo mismo que les voy a dar yo, o tal vez mejor. Entonces tenemos que aprender a utilizar estas nuevas herramientas.''}
\end{quote}

La dinámica de trabajo consiste en que los estudiantes lean, escuchen o investiguen los temas antes de cada clase, de modo que el tiempo de clase se dedique a la discusión y aplicación práctica de esos conceptos.

\begin{important}[]
\textbf{Regla de uso de la inteligencia artificial:} los estudiantes pueden consultar herramientas como ChatGPT para resolver un caso o problema planteado, pero la respuesta obtenida debe pasar siempre por un filtro propio (``¿Ustedes actuarían de esa manera? ¿Sería lo correcto para ustedes?''). Lo que se entrega no es la respuesta de la inteligencia artificial, sino la conclusión personal construida a partir de ella.
\end{important}

Como técnica de estudio, la docente pidió que después de cada clase se envíe un breve resumen de cuatro palabras sobre el tema tratado, explicando el fundamento pedagógico del ejercicio:

\begin{quote}
\textit{``Cuando ustedes cierran un tema y lo escriben, esto de escribir genera en la memoria algo diferente que si solamente lo digo y lo hablo. Cuando yo resumo, me queda en la memoria.''}
\end{quote}

\subsection{Valor de la experiencia profesional temprana}

\begin{note}
Varios estudiantes consultaron sobre la conveniencia de trabajar mientras cursan, aun sin necesidad económica. La docente recomendó siempre trabajar de lo que sea, aunque no haga falta por dinero, porque ``te da una experiencia de la calle que no te la da la facultad''. Relató que ella misma trabajó gratis en un estudio jurídico y luego en una escribanía durante varios años al inicio de su carrera, en un contexto económico adverso, y que no se arrepiente: ``la experiencia que logré no tiene precio''.
\end{note}

Por último, se planteó que la relación entre docente y estudiantes, en esta etapa final de la carrera, no debe regirse por una jerarquía rígida, sino por un vínculo de consenso y aprendizaje mutuo.

\section{El mundo del trabajo en transformación}

La docente introdujo el eje de la clase señalando que el derecho no puede permanecer ajeno a una transformación del mundo del trabajo que viene anunciándose desde hace años:

\begin{quote}
\textit{``La globalización se produjo, la pandemia se produjo, y esa transformación entró con todo.''}
\end{quote}

Explicó que la pandemia no originó el cambio, sino que aceleró un proceso de transformación que ya estaba en marcha, modificando de manera estructural la forma en que las personas trabajan y se vinculan profesionalmente.

\begin{example}{La nueva ciudadanía del trabajo remoto}
La docente relató dos encuentros personales para ilustrar el fenómeno: un argentino que se mudó a vivir a Australia y seguía manteniendo, a la distancia, sus costumbres y su hinchada de fútbol como si nunca se hubiera ido; y una persona cuyo trabajo, por ser enteramente online, le permitía ejercerlo desde cualquier lugar del mundo, alternando temporadas en distintos países.

La conclusión extraída de estos casos fue que hoy ``no importa dónde vivís'', porque el trabajo se desprendió del territorio. Esa desterritorialización exige repensar también el derecho que regula la actividad económica y las profesiones.
\end{example}

A partir de estos ejemplos, la docente advirtió sobre la obsolescencia de ciertas expectativas laborales: profesiones y modos de trabajo que hoy parecen sólidos podrían dejar de existir en pocos años, por lo cual instó al grupo a pensar su formación con una mirada flexible y de mediano plazo.

\section{La reforma proyectada de la ley societaria}

La docente adelantó que el Congreso tiene en agenda un proyecto de reforma integral de la Ley General de Sociedades, y adelantó su valoración general: comparte algunos de sus lineamientos y discrepa con otros, pero destacó como mérito central que el proyecto no se limita a corregir la ley vigente, sino que propone un cambio de paradigma.

\begin{quote}
\textit{``No copia lo anterior. Deja totalmente de lado todo lo anterior y ve el mundo de hoy como está.''}
\end{quote}

Para explicar la diferencia entre una reforma ``cosmética'' y una reforma estructural, se recurrió a una metáfora que atravesó buena parte de la clase:

\begin{quote}
\textit{``Le cambio el vestido, pero el cuerpo sigue igual. Le pongo un vestido más de joven, pero ya tiene un cuerpo viejo: por más que me ponga una minifalda, sigo siendo vieja.''}
\end{quote}

\begin{example}{El antecedente del proyecto de reforma de 2018}
La docente relató que, como responsable del Instituto de Derecho Comercial del colegio de escribanos, impulsó en el año 2018 el debate del primer proyecto de reforma de la ley de sociedades, lo cual le trajo fuertes resistencias internas dentro de la institución.

Calificó a ese proyecto de 2018 como ``muy bueno'' desde el punto de vista doctrinario, pero limitado, porque modificaba aspectos puntuales sin transformar la estructura de fondo: seguía siendo el mismo ``cuerpo'' con un ``vestido'' distinto. El proyecto actual, en cambio, no reforma sobre lo existente sino que parte de una base nueva.
\end{example}

Como antecedente de lo que implica un cambio normativo estructural, se recordó la experiencia profesional con la sanción del Código Civil y Comercial: de un día para el otro dejó de poder aplicarse el código anterior, lo que obligó a un reentrenamiento intensivo de todos los operadores jurídicos.

\begin{quote}
\textit{``Yo me recibí con el código viejo, toda mi vida trabajé con el código viejo, y de un día para otro tuvimos que aplicar el nuevo. Tuvimos que hacer un curso, y otro curso, y otro curso, para poder empezar a trabajar.''}
\end{quote}

Se cerró la unidad señalando que, aun cuando el nuevo proyecto de ley de sociedades no llegue a sancionarse, ya está operando un cambio de paradigma en la manera de mirar el derecho societario.

% ---------------- CORNELL CAPÍTULO 1 ----------------
\section*{Cuadro Cornell — Presentación y contexto del cambio}
\addcontentsline{toc}{section}{Cuadro Cornell — Presentación y contexto del cambio}

\begin{table}[H]
\centering
\begin{tabularx}{\textwidth}{
    >{\raggedright\arraybackslash}p{0.32\textwidth}
    >{\raggedright\arraybackslash}X
}
\toprule
\textbf{Preguntas / palabras clave} &
\textbf{Notas / respuestas} \\
\midrule

\textbf{¿Por qué ningún curso de la cátedra es igual al anterior?} &
Porque la profundidad y el enfoque de cada clase se adaptan a lo que el grupo exige y está dispuesto a dar. La base del programa es fija, pero lo construido sobre ella depende del compromiso de cada curso en particular. \\

\textbf{¿Qué rol cumple la inteligencia artificial en la cátedra?} &
Es una herramienta de consulta legítima, pero la respuesta que ofrece debe pasar siempre por un criterio propio antes de presentarse como trabajo: se entrega la conclusión personal, no la respuesta textual de la herramienta. \\

\textbf{¿Por qué el derecho societario está en transición?} &
Porque la globalización y la pandemia aceleraron una transformación del mundo del trabajo que ya estaba en marcha (desterritorialización del trabajo, digitalización), y el derecho debe adaptarse a esa nueva realidad económica. \\

\textbf{¿Qué diferencia hay entre el proyecto de reforma societaria de 2018 y el actual?} &
El de 2018 modificaba aspectos puntuales de la ley vigente sin cambiar su estructura de fondo (``cambiar el vestido''). El proyecto actual parte de cero y propone un cambio de paradigma (``cambiar el cuerpo''). \\

\midrule

\multicolumn{2}{p{\dimexpr\textwidth-2\tabcolsep\relax}}{
\textbf{Resumen:} La cátedra se organiza como un proceso adaptable al grupo y apoyado en el uso crítico de la inteligencia artificial. Sobre ese marco pedagógico se introduce el eje central del curso: el derecho societario atraviesa una transición motivada por la transformación acelerada del mundo del trabajo, que impulsa una reforma proyectada de la Ley General de Sociedades concebida como un cambio de paradigma y no como una simple actualización cosmética.
} \\

\bottomrule
\end{tabularx}
\end{table}

% ============================================================
% CAPÍTULO 2
% ============================================================
\chapter[Personas y fuentes del derecho]{Personas y fuentes del derecho}

\section{Persona humana y persona jurídica}

Se ubicó como punto de partida conceptual la distinción entre persona humana y persona jurídica, subrayando su relevancia práctica para el ejercicio profesional: la mayoría de las contrapartes con las que un abogado, contador o escribano trabajará —ya sea para contratar, demandar o defender— serán personas jurídicas, y no personas humanas.

\begin{important}
Antes de contratar, demandar, impugnar una decisión o verificar una nulidad frente a una persona jurídica, el profesional debe verificar: si quien actúa en su representación tenía facultades suficientes; si esa representación era correcta; y si la toma de decisión (asamblea o reunión de socios) cumplió con los pasos exigidos según el objeto social.
\end{important}

Se adelantó que el concepto de ``objeto social'' está en plena transformación: bajo el régimen vigente, la responsabilidad se organiza en gran medida en torno al objeto social, mientras que el proyecto de reforma societaria lo desplaza a un segundo plano, dando centralidad a la actividad efectivamente desarrollada.

\section{El Código Civil y Comercial y la autonomía de la voluntad}

Se explicó que el Código Civil y Comercial unificó los conceptos societarios, comerciales y civiles en un único cuerpo normativo. Como consecuencia, la figura del derecho comercial como rama jurídica autónoma dejó de existir formalmente, aunque en la práctica los tribunales con competencia comercial se mantienen.

\begin{quote}
\textit{``El derecho comercial como forma independiente no existe más. Sin embargo, en la práctica todavía seguimos teniendo jueces en lo comercial.''}
\end{quote}

Se valoró positivamente conservar la mirada comercial junto a la civil para el análisis de un mismo caso, por entender que cada perspectiva aporta matices distintos y complementarios.

\subsection{Autonomía de la voluntad: libertades y restricciones}

Se definió a la autonomía de la voluntad como eje central de la mirada del derecho expuesta en la clase, y se explicó que el Código Civil y Comercial es asimétrico en el modo en que la garantiza: amplia libertad en materia de decisiones personalísimas (por ejemplo, sobre el propio cuerpo y la salud) y, en cambio, fuertes restricciones en materia patrimonial y sucesoria.

\begin{definition}{La porción legítima hereditaria}
El Código Civil y Comercial reserva a los herederos forzosos (descendientes, ascendientes y cónyuge) una porción del patrimonio de la que el titular no puede disponer libremente por testamento, conocida como legítima hereditaria.

\textit{Fuente: Código Civil y Comercial de la Nación, régimen de la porción legítima (arts. 2444 y siguientes).}
\end{definition}

Se adelantó una postura personal crítica frente a esta institución, autodefinida como ``antilegitimaria'':

\begin{quote}
\textit{``Yo pienso que con mi patrimonio puedo hacer lo que se me antoje. Si tuve la suerte de tener hijos maravillosos, van a recibir lo mío porque yo quiero. Pero si tuviera hijos que no me dieron bola, no tengo por qué tener que inventar instrumentos jurídicos para poder dejarle a la persona que efectivamente me cuidó.''}
\end{quote}

Se precisó que se trata de una opinión personal, no de una posición unánime, sostenida incluso trabajando profesionalmente en la consultoría de empresas familiares, donde la planificación patrimonial y sucesoria es moneda corriente.

\section{La pirámide de Kelsen y los valores del intérprete}

Para explicar por qué cada profesional puede llegar a soluciones distintas frente a un mismo caso, se recurrió a la pirámide de Kelsen, no ya como esquema formal de jerarquía normativa, sino como metáfora del sistema de valores personales de quien interpreta el derecho.

\begin{quote}
\textit{``En la punta de la pirámide van a estar sus valores, los valores con los que cada uno de ustedes labora. Si para ustedes lo máximo es la Constitución Nacional, es una cosa; para mí lo máximo es el derecho natural, porque yo creo en Dios.''}
\end{quote}

\begin{example}{El semáforo apagado}
Se ejemplificó la diferencia entre actuar por temor a la sanción y actuar por convicción con una situación cotidiana: frente a un semáforo apagado, la docente se detiene igual, sin que nadie se lo exija, porque entiende que es lo correcto —tanto en el tránsito como al conducir su vida profesional—. Ese es el concepto de responsabilidad que, a su juicio, todavía falta consolidar en el ejercicio societario: no cumplir la norma para evitar una sanción, sino cumplirla porque corresponde.
\end{example}

Se explicó que en el ámbito universitario y profesional coexisten distintas escuelas de interpretación del derecho societario, cada una con su propia jerarquía de valores; se mencionaron, a título ilustrativo y sin ánimo de agotar el panorama doctrinario, las posiciones más conservadoras —identificadas con juristas como Ricardo Nissen— frente a lecturas más flexibles —identificadas con autores como Rafael Manóvil—, sosteniéndose que ambas son igualmente válidas como puntos de partida.

\begin{important}
No existe una única forma ``correcta'' de mirar el derecho societario. La docente se comprometió a exponer su propia mirada sin desconocer la validez de otros criterios, para no sesgar la información que reciben los estudiantes.
\end{important}

También se citó al jurista cordobés Efraín Hugo Richard —fallecido recientemente— como referencia doctrinaria central para pensar la relación entre libertad y responsabilidad societaria, resumida en la fórmula adoptada como eje de la unidad siguiente:

\begin{quote}
\textit{``A mayor libertad, mayor responsabilidad.''}
\end{quote}

% ---------------- CORNELL CAPÍTULO 2 ----------------
\section*{Cuadro Cornell — Personas y fuentes del derecho}
\addcontentsline{toc}{section}{Cuadro Cornell — Personas y fuentes del derecho}

\begin{table}[H]
\centering
\begin{tabularx}{\textwidth}{
    >{\raggedright\arraybackslash}p{0.32\textwidth}
    >{\raggedright\arraybackslash}X
}
\toprule
\textbf{Preguntas / palabras clave} &
\textbf{Notas / respuestas} \\
\midrule

\textbf{¿Qué es la autonomía de la voluntad y dónde encuentra sus límites en el Código Civil y Comercial?} &
Es el principio que habilita a las personas a organizar libremente su vida jurídica. El Código Civil y Comercial la garantiza ampliamente en materia de decisiones personalísimas (salud, cuerpo), pero la restringe en materia patrimonial y sucesoria, por ejemplo a través de la porción legítima hereditaria. \\

\textbf{¿Qué representa la pirámide de Kelsen en esta clase?} &
No se usa como esquema formal de jerarquía normativa, sino como metáfora del sistema de valores personales de quien interpreta el derecho: en la cima de la pirámide de cada persona está lo que considera su fuente de sentido última (la ley, la Constitución, el derecho natural, etc.). \\

\midrule

\multicolumn{2}{p{\dimexpr\textwidth-2\tabcolsep\relax}}{
\textbf{Resumen:} La unificación del derecho civil y comercial en el Código Civil y Comercial reorganizó las fuentes del derecho societario en torno a la autonomía de la voluntad, garantizada de modo asimétrico según se trate de decisiones personalísimas o patrimoniales. La interpretación de esas fuentes depende, a su vez, del sistema de valores propio de cada operador jurídico, lo que explica que frente a un mismo caso coexistan lecturas doctrinarias distintas e igualmente válidas.
} \\

\bottomrule
\end{tabularx}
\end{table}

% ============================================================
% CAPÍTULO 3
% ============================================================
\chapter[El sistema de tipicidad clásica]{El sistema de tipicidad clásica}

\section{Los tipos societarios previstos por la ley}

Se introdujo el principio de tipicidad que estructura a la Ley General de Sociedades N.\textordmasculine~19.550: para que una organización sea considerada sociedad en los términos de esa ley, debe adoptar uno de los tipos expresamente previstos.

\begin{formula}{Principio de tipicidad societaria}
La Ley General de Sociedades N.\textordmasculine~19.550 exige, para constituir una sociedad regularmente encuadrada en su régimen, la adopción de uno de los tipos previstos por la propia ley. Toda organización que no adopte alguno de esos tipos queda fuera del sistema de tipicidad, con un régimen jurídico distinto (sociedad simple o residual).

\textit{Fuente: Ley General de Sociedades N.\textordmasculine~19.550, régimen de tipicidad.}
\end{formula}

Los tipos previstos, ordenados de mayor a menor complejidad estructural, son:

\begin{itemize}
    \item Sociedad Anónima (S.A.)
    \item Sociedad en Comandita por Acciones
    \item Sociedad de Responsabilidad Limitada (S.R.L.)
    \item Sociedad en Comandita Simple
    \item Sociedad Colectiva
    \item Sociedad de Capital e Industria
\end{itemize}

Se adelantó el criterio que recorrería el resto de la clase: contar con un tipo previsto por la ley no equivale a que ese tipo sea el adecuado para cada caso concreto. Elegir el tipo correcto —y no automáticamente el más conocido o prestigioso— es la tarea central del asesoramiento profesional.

\section{Sociedad Anónima y Sociedad de Responsabilidad Limitada}

Se explicó el origen histórico de la S.R.L. como un tipo intermedio, creado por el legislador de 1972 con un perfil más personalista que la S.A., ante la constatación de que muchos pequeños comerciantes no necesitaban la estructura corporativa de una sociedad anónima:

\begin{quote}
\textit{``La gente no se da cuenta de que es lo mismo ser gerente de una S.R.L. que presidente de una S.A. Y entonces se tiran a hacer la S.A. para un kiosco, cuando ese ropaje jurídico no le sirve.''}
\end{quote}

\begin{example}{El kiosco constituido como S.A.}
Se utilizó el caso de un comerciante que constituye una sociedad anónima para un negocio familiar pequeño (``un kiosco'') como paradigma del error de asesoramiento. La S.A. exige reuniones de directorio cada tres meses, aun cuando los socios sean, por ejemplo, un matrimonio que se ve todos los días.

La consecuencia no es solamente el costo y la carga administrativa de sostener el tipo: la falta de libros societarios actualizados y de actas de directorio puede convertirse en una herramienta para impugnar decisiones sociales o para habilitar reclamos de terceros —por ejemplo, un reclamo laboral de quien trabajó sin estar registrado, que puede valerse de la falta de libros en regla para reforzar su posición—.
\end{example}

A modo de síntesis, se contrapusieron ambos tipos en sus principales exigencias:

\begin{table}[H]
\centering
\caption{Sociedad Anónima frente a Sociedad de Responsabilidad Limitada}
\begin{tabularx}{\textwidth}{>{\raggedright\arraybackslash}p{0.22\textwidth}XX}
\toprule
\textbf{Aspecto} & \textbf{Sociedad Anónima (S.A.)} & \textbf{Sociedad de Resp. Limitada (S.R.L.)} \\
\midrule
Órgano de administración & Directorio, renovación máxima cada 3 años & Gerencia, por tiempo indeterminado \\
Reuniones obligatorias & Directorio cada 3 meses como máximo & Reunión de socios anual \\
Presentaciones ante IGJ & Mayor carga de presentaciones periódicas & Sin presentación de actas de gerencia \\
Perfil & Corporativo, pensado para grandes estructuras & Personalista, pensado para grupos reducidos \\
\bottomrule
\end{tabularx}
\end{table}

\begin{important}
El tipo societario no es un tema meramente formal, sino una decisión de asesoramiento profesional con consecuencias patrimoniales, tributarias y hasta laborales: hay que elegir la estructura acorde a la realidad económica del negocio, no a la que parece más prestigiosa.
\end{important}

\section{Sociedades en extinción}

Se describieron tres tipos societarios en franco desuso: la sociedad en comandita simple, la sociedad colectiva y la sociedad de capital e industria. Se señaló que, en toda la trayectoria profesional relatada, nunca se constituyó una sociedad de estos tipos desde cero; solo se trabajó con sociedades preexistentes de este tipo para transformarlas en otras.

Sin embargo, se destacó que el perfil marcadamente personalista de la sociedad colectiva podría ser aprovechado hoy para estructurar determinados contratos asociativos que no encuentran un molde adecuado en los tipos más utilizados.

\begin{example}{La sociedad de capital e industria y el startup actual}
Se equiparó la sociedad de capital e industria —donde un socio aporta el trabajo y el otro el capital— con la lógica de financiamiento de un startup actual: una persona aporta la idea y el trabajo, y otra el dinero. El problema es que ese tipo societario, tal como está regulado, ya no refleja adecuadamente esa relación.

En un startup real, quien aporta la idea no puede quedarse con el 50\% del capital social porque no está aportando capital, pero tampoco quien pone el dinero puede quedarse con la idea. Forzar esa relación dentro de una sociedad de capital e industria obliga, en la práctica, a distorsionar en el instrumento constitutivo los aportes reales de cada parte, y expone además a quien aportó la idea a perder poder de decisión si más adelante se diluye su participación en una futura ronda de inversión.
\end{example}

% ---------------- CORNELL CAPÍTULO 3 ----------------
\section*{Cuadro Cornell — El sistema de tipicidad clásica}
\addcontentsline{toc}{section}{Cuadro Cornell — El sistema de tipicidad clásica}

\begin{table}[H]
\centering
\begin{tabularx}{\textwidth}{
    >{\raggedright\arraybackslash}p{0.32\textwidth}
    >{\raggedright\arraybackslash}X
}
\toprule
\textbf{Preguntas / palabras clave} &
\textbf{Notas / respuestas} \\
\midrule

\textbf{¿Qué significa el principio de tipicidad societaria?} &
Que, para quedar encuadrada en el régimen general de la Ley de Sociedades, una sociedad debe adoptar uno de los tipos previstos por la ley (S.A., S.R.L., comandita simple, comandita por acciones, colectiva o capital e industria). \\

\textbf{¿Por qué una S.A. puede ser un ``traje'' equivocado para un negocio pequeño?} &
Porque impone cargas estructurales —reuniones de directorio periódicas, libros societarios, presentaciones ante el organismo de control— que generan costos y riesgos legales (impugnaciones, reclamos laborales) innecesarios para una estructura familiar o de pocos socios, donde una S.R.L. suele ser más adecuada. \\

\midrule

\multicolumn{2}{p{\dimexpr\textwidth-2\tabcolsep\relax}}{
\textbf{Resumen:} El sistema clásico de tipicidad obliga a adoptar uno de los tipos societarios previstos por la Ley General de Sociedades. Entre los tipos vigentes, la Sociedad Anónima y la Sociedad de Responsabilidad Limitada son los de mayor uso práctico, pero su aplicación sin criterio profesional —como constituir una S.A. para un negocio familiar pequeño— genera cargas y riesgos innecesarios. Otros tipos, como la comandita simple, la colectiva y la de capital e industria, se encuentran hoy en franco desuso, aunque su lógica personalista conserva utilidad para estructurar ciertos contratos asociativos.
} \\

\bottomrule
\end{tabularx}
\end{table}

% ============================================================
% CAPÍTULO 4
% ============================================================
\chapter[La crisis del tipo y la sociedad unipersonal]{La crisis del tipo y la sociedad unipersonal}

\section{La sociedad simple (ex Sección IV)}

Se explicó que, con la sanción del Código Civil y Comercial, las sociedades civiles —que hasta entonces convivían con las sociedades comerciales tipificadas— desaparecieron como categoría autónoma. En su lugar, ganó protagonismo una figura que ya existía en la Ley de Sociedades pero con otro perfil: la sociedad simple, regulada en la Sección IV de la ley.

\begin{quote}
\textit{``Entra la niña bonita: la sociedad simple. Esa era antes la sociedad de hecho irregular, que no era querida por el legislador de 1972, y ahora empieza a tener una preponderancia que antes no tenía.''}
\end{quote}

\begin{definition}{La sociedad simple o residual}
Las sociedades que no se constituyen según alguno de los tipos previstos por la Ley General de Sociedades, que omiten requisitos esenciales o que incumplen con las formalidades exigidas, quedan regidas por el régimen de la Sección IV de la ley, conocido en la práctica como régimen de la ``sociedad simple''. A diferencia del sistema anterior a la reforma, este régimen ya no sanciona la irregularidad con la nulidad, sino que reconoce a estas sociedades personalidad jurídica plena, con reglas propias de responsabilidad y oponibilidad.

\textit{Fuente: Ley General de Sociedades N.\textordmasculine~19.550, Sección IV, arts. 21 a 26.}
\end{definition}

La ventaja central de este régimen es que permite a los socios pactar libremente su organización interna sin sujetarse a la rigidez de un tipo previsto, a costa de asumir una responsabilidad distinta frente a terceros. Se retomó aquí la fórmula de Hugo Richard citada en la unidad anterior:

\begin{quote}
\textit{``Esto es lo que me da la autonomía de la voluntad. Yo me pongo el vestido que quiero con la sociedad simple, pero a cambio asumo otra responsabilidad. Libertad con responsabilidad: ese es el eje.''}
\end{quote}

\section{La Sociedad Anónima Unipersonal (S.A.U.)}

Se narró el origen histórico de la S.A.U. como resultado de un largo debate doctrinario: en un congreso de derecho societario se discutieron más de veinte proyectos distintos para regular la posibilidad de que una sola persona constituyera una sociedad, cada uno con su propio nombre.

El proyecto original —tal como fue presentado en el debate parlamentario del Código Civil y Comercial— apuntaba a dar una herramienta ágil al pequeño emprendedor que no quería (o no podía) sumar un socio formal solo para cumplir con el requisito de pluralidad de socios.

\begin{example}{El ``socio de paja''}
Se relató una situación habitual en los estudios jurídicos previa a la existencia de la sociedad unipersonal: un cliente que necesita constituir una sociedad pero no tiene con quién asociarse, y termina ofreciendo la calidad de socio a su cónyuge, a un hermano con el que está distanciado, o a un familiar de edad muy avanzada, solo para cumplir el requisito formal de pluralidad de socios.

Ese ``socio de paja'', aunque tenga una participación mínima —por ejemplo, del 5\%—, conserva igualmente derechos societarios como el de convocar a asamblea, lo cual puede generar conflictos societarios no queridos por el socio mayoritario.
\end{example}

Sin embargo, cuando el proyecto pasó por el Poder Ejecutivo antes de su sanción definitiva, se modificó su perfil original y se lo redujo a la actual Sociedad Anónima Unipersonal, con requisitos que desvirtuaron el objetivo inicial de asistir al pequeño emprendedor:

\begin{quote}
\textit{``Una mano limpió eso y generó la Sociedad Anónima Unipersonal, que desbarató el objetivo que tenía, porque la unipersonalidad estaba pensada en mayor medida para el pequeño emprendedor, para el pequeño comerciante que no quiere armar una sociedad.''}
\end{quote}

\begin{formula}{Requisitos de la Sociedad Anónima Unipersonal}
La constitución de una S.A.U. exige: (i) un capital social mínimo de \$30.000.000, actualizado por el Decreto N.\textordmasculine~209/2024 (que modificó el art. 186 de la Ley General de Sociedades N.\textordmasculine~19.550, elevando el monto que regía desde 2012); (ii) la integración del 100\% del capital social en el mismo acto constitutivo, a diferencia de la sociedad anónima pluripersonal, donde la ley permite integrar un 25\% al inicio y el saldo dentro de los dos años; y (iii) sindicatura obligatoria, con el consecuente sometimiento al régimen de fiscalización estatal permanente del art. 299 de la ley.

\textit{Fuente: Ley General de Sociedades N.\textordmasculine~19.550, arts. 186 y 299; Decreto N.\textordmasculine~209/2024.}
\end{formula}

Se subrayó especialmente la exigencia de sindicatura obligatoria como el mayor obstáculo para el pequeño emprendedor, que no está en condiciones de afrontar ese costo adicional ni desea quedar sujeto a un régimen de control estatal reforzado:

\begin{quote}
\textit{``Si yo soy un pequeño comerciante que quiere adoptar una S.A.U., no quiero tener control estatal permanente y no puedo pagar una sindicatura porque tengo que pagar un montón de plata. Entonces se desvirtúa el objetivo de la S.A.U. ¿Y para quién quedó? Para las sociedades vehículo, para las sociedades extranjeras que vienen a constituir una S.A.U. porque no quieren tener un socio externo. Y ahí está perfecto.''}
\end{quote}

\begin{important}
La S.A.U., pensada originalmente como una herramienta de simplificación para el pequeño emprendedor, terminó siendo utilizada casi exclusivamente por grandes grupos económicos y sociedades extranjeras que necesitan una sociedad vehículo sin socios externos. El pequeño comerciante, en la práctica, sigue necesitando otras alternativas —como la sociedad simple o la S.A.S.— para evitar recurrir al ``socio de paja''.
\end{important}

\section{Suscripción e integración del capital social}

Se distinguieron con precisión dos conceptos que suelen confundirse en la práctica:

\begin{definition}{Suscripción vs. integración}
\textbf{Suscripción:} es el monto de capital que el socio se compromete a aportar a la sociedad.

\textbf{Integración:} es el monto que el socio efectivamente aporta a la sociedad en un momento dado.
\end{definition}

Se explicó que la posibilidad de suscribir un capital e integrar solo una parte de él en el acto constitutivo es una excepción reservada exclusivamente a los aportes en dinero en efectivo:

\begin{quote}
\textit{``En el único caso en que puedo suscribir cien e integrar menos cantidad es cuando aporto dinero. Si aporto otro tipo de bienes, tengo que suscribir e integrar la totalidad. Si digo que aporto dólares, tengo que suscribir e integrar la totalidad. Si digo que aporto un inmueble, suscribo e integro la totalidad.''}
\end{quote}

Para los aportes dinerarios, la ley permite integrar un porcentaje inicial (25\% en la sociedad anónima pluripersonal) y completar el saldo dentro de un plazo de hasta dos años, según lo que determine el directorio.

\begin{example}{La venta del departamento}
Se propuso el caso de un cliente que quiere constituir una sociedad con un capital social elevado, pero no dispone en el momento de la totalidad del dinero. La solución práctica es suscribir el total del capital previsto e integrar solo una parte —por ejemplo, un 25\%— al momento de la constitución, comprometiéndose a integrar el saldo dentro de los días siguientes, cuando cobre, por ejemplo, el producido de la venta de un inmueble. La ley da así un margen de planificación financiera para el cliente, siempre que se trate de aportes dinerarios.
\end{example}

% ---------------- CORNELL CAPÍTULO 4 ----------------
\section*{Cuadro Cornell — La crisis del tipo y la sociedad unipersonal}
\addcontentsline{toc}{section}{Cuadro Cornell — La crisis del tipo y la sociedad unipersonal}

\begin{table}[H]
\centering
\begin{tabularx}{\textwidth}{
    >{\raggedright\arraybackslash}p{0.32\textwidth}
    >{\raggedright\arraybackslash}X
}
\toprule
\textbf{Preguntas / palabras clave} &
\textbf{Notas / respuestas} \\
\midrule

\textbf{¿Qué es la sociedad simple (ex Sección IV) y qué ventaja ofrece?} &
Es el régimen que rige a las sociedades que no adoptan ningún tipo previsto por la ley. A cambio de mayor libertad de organización interna, impone un régimen de responsabilidad distinto: ``a mayor libertad, mayor responsabilidad''. \\

\textbf{¿Cuáles son los requisitos centrales de la Sociedad Anónima Unipersonal (S.A.U.)?} &
Capital social mínimo de \$30.000.000 (Decreto 209/2024), integración del 100\% del capital en el acto constitutivo y sindicatura obligatoria, con sujeción al régimen de fiscalización estatal permanente del art. 299 de la Ley General de Sociedades. \\

\textbf{¿Por qué se dice que la S.A.U. ``desbarató'' su objetivo original?} &
Porque fue pensada para asistir al pequeño emprendedor que no quería recurrir a un ``socio de paja'', pero terminó siendo una herramienta usada principalmente por grandes sociedades vehículo y sociedades extranjeras, dado el costo del capital mínimo y de la sindicatura obligatoria. \\

\textbf{¿Cuál es la diferencia entre suscripción e integración del capital?} &
La suscripción es el monto que el socio se compromete a aportar; la integración es el monto que efectivamente aporta. Solo en los aportes dinerarios la ley permite integrar una parte al constituir la sociedad (25\%) y el resto dentro de dos años; los aportes en especie deben integrarse en su totalidad desde el inicio. \\

\midrule

\multicolumn{2}{p{\dimexpr\textwidth-2\tabcolsep\relax}}{
\textbf{Resumen:} Frente a la rigidez del sistema de tipicidad, surgen dos vías de flexibilización con resultados prácticos muy distintos. La sociedad simple prescinde de un tipo previsto a cambio de un régimen de responsabilidad distinto, mientras que la Sociedad Anónima Unipersonal —pese a haber sido pensada para el pequeño emprendedor— terminó exigiendo un capital mínimo elevado y sindicatura obligatoria, quedando reservada en la práctica a grandes sociedades vehículo. Se agrega además la distinción técnica entre suscripción e integración del capital social, relevante para planificar los aportes de los socios.
} \\

\bottomrule
\end{tabularx}
\end{table}

% ============================================================
% CAPÍTULO 5
% ============================================================
\chapter[La sociedad por acciones simplificada]{La sociedad por acciones simplificada}

\section{La Ley de Apoyo al Capital Emprendedor y el nacimiento de la S.A.S.}

Se ubicó el nacimiento de la Sociedad por Acciones Simplificada (S.A.S.) dentro de la Ley de Apoyo al Capital Emprendedor, explicándose la técnica legislativa utilizada: incorporar figuras societarias novedosas a través de una ley específica, por fuera de la Ley General de Sociedades, cuando modificar esta última resulta políticamente más difícil.

\begin{definition}{La Sociedad por Acciones Simplificada (S.A.S.)}
La Ley N.\textordmasculine~27.349, de Apoyo al Capital Emprendedor, creó la Sociedad por Acciones Simplificada como un nuevo tipo societario por fuera de la Ley General de Sociedades N.\textordmasculine~19.550, con constitución íntegramente digital, objeto amplio, mayor libertad para pactar la representación y la administración, y aplicación supletoria en primer lugar del Código Civil y Comercial y, en segundo lugar, de la Ley General de Sociedades ante lagunas normativas.

\textit{Fuente: Ley N.\textordmasculine~27.349, de Apoyo al Capital Emprendedor.}
\end{definition}

Se destacaron como diferencias centrales de la S.A.S. frente a los tipos tradicionales:

\begin{itemize}
    \item Permite la unipersonalidad sin las exigencias de capital mínimo elevado ni sindicatura obligatoria que pesan sobre la S.A.U.
    \item Su constitución puede tramitarse en soporte íntegramente digital, sin necesidad de instrumento en papel.
    \item El objeto social puede redactarse en forma amplia, sin la exigencia de precisión y determinación que tradicionalmente exige la Ley General de Sociedades.
    \item Permite diseñar formas de representación y de administración más flexibles, incluyendo mandatos por tiempo indeterminado.
\end{itemize}

\begin{note}
Se consultó con qué criterio se decide, ante una laguna normativa, si aplicar el Código Civil y Comercial o la Ley General de Sociedades a una S.A.S. Se explicó que el orden de prelación es inverso al de las sociedades típicas: frente a una laguna en la S.A.S., corresponde recurrir primero al Código Civil y Comercial y, solo si este tampoco da respuesta, a la Ley General de Sociedades.
\end{note}

Se advirtió que la S.A.S. generó resistencia en un sector de la doctrina societaria más conservadora, para la cual la flexibilidad del nuevo régimen ponía en riesgo la seguridad jurídica que aportaba el sistema clásico de tipicidad. Se relató, a modo de anécdota, la reacción de colegas de mayor trayectoria ante una charla sobre esta figura: le señalaron que estaba ``destruyendo la notarial'' al restarle protagonismo a la escritura pública en la constitución de sociedades.

A modo de síntesis comparativa entre los tres tipos societarios de mayor uso práctico:

\begin{table}[H]
\centering
\caption{S.A., S.R.L. y S.A.S.: comparación de exigencias centrales}
\begin{tabularx}{\textwidth}{>{\raggedright\arraybackslash}p{0.22\textwidth}XXX}
\toprule
\textbf{Aspecto} & \textbf{S.A.} & \textbf{S.R.L.} & \textbf{S.A.S.} \\
\midrule
Capital mínimo & \$30.000.000 (Dto. 209/2024) & Sin mínimo legal & Sin mínimo legal elevado \\
Constitución digital & No, requiere escritura o instrumento privado & No, requiere instrumento privado & Sí, íntegramente digital \\
Sindicatura obligatoria & Solo en casos del art. 299 LGS & No & No \\
Norma supletoria & Ley General de Sociedades & Ley General de Sociedades & Primero CCyC, luego LGS \\
\bottomrule
\end{tabularx}
\end{table}

\section{La forma frente al fondo: escritura pública y firma digital}

Se cerró la clase con una reflexión sobre el valor jurídico de la forma, en diálogo directo con la resistencia doctrinaria descripta en el apartado anterior. Se sostuvo que el soporte de un acto jurídico —papel, disco rígido o nube— no determina su validez ni su eficacia, sino el cumplimiento efectivo de la función que ese acto está llamado a cumplir:

\begin{quote}
\textit{``¿Qué me importa la escritura pública? No importa que sea el vehículo necesario para que el empresario pueda operar. La forma no hace el fondo. Es lo mismo que esté en un papel, que esté en un disco rígido, que esté en una nube: lo importante es la tarea que uno realiza.''}
\end{quote}

En esa misma línea, se expresó una expectativa favorable respecto de la implementación de un protocolo notarial digital, argumentando que la certificación digital de una firma cumple la misma función de dar fe pública que el protocolo en papel, sin necesidad de sostener las cargas físicas tradicionales del oficio notarial.

\begin{note}
Se mencionó, como antecedente de la técnica legislativa de incorporar institutos nuevos a través de leyes especiales, la Ley N.\textordmasculine~24.441, que introdujo en el ordenamiento argentino la figura del fideicomiso sin necesidad de reformar el Código Civil entonces vigente.
\end{note}

% ---------------- CORNELL CAPÍTULO 5 ----------------
\section*{Cuadro Cornell — La sociedad por acciones simplificada}
\addcontentsline{toc}{section}{Cuadro Cornell — La sociedad por acciones simplificada}

\begin{table}[H]
\centering
\begin{tabularx}{\textwidth}{
    >{\raggedright\arraybackslash}p{0.32\textwidth}
    >{\raggedright\arraybackslash}X
}
\toprule
\textbf{Preguntas / palabras clave} &
\textbf{Notas / respuestas} \\
\midrule

\textbf{¿Qué es la Sociedad por Acciones Simplificada (S.A.S.) y qué la distingue?} &
Es un tipo societario creado por la Ley de Apoyo al Capital Emprendedor (Ley 27.349), por fuera de la Ley General de Sociedades. Permite constitución íntegramente digital, objeto amplio, mayor flexibilidad en la representación, y aplica supletoriamente primero el Código Civil y Comercial y luego la Ley General de Sociedades. \\

\textbf{¿Qué significa la frase ``la forma no hace el fondo''?} &
Que el soporte de un acto jurídico —papel, disco rígido o nube— no determina su validez, sino el cumplimiento efectivo de la función que ese acto está llamado a cumplir. Es el argumento central a favor de la digitalización de la función notarial. \\

\midrule

\multicolumn{2}{p{\dimexpr\textwidth-2\tabcolsep\relax}}{
\textbf{Resumen:} La Sociedad por Acciones Simplificada, creada por la Ley de Apoyo al Capital Emprendedor, ofrece hoy la alternativa más flexible y digital para el emprendedor real, frente a las limitaciones prácticas de la S.A. y la S.A.U. El cierre conceptual de la clase —``la forma no hace el fondo''— conecta este tema con una reflexión más amplia sobre el futuro de las profesiones jurídicas: lo que importa no es el soporte de un acto, sino la función que ese acto efectivamente cumple.
} \\

\bottomrule
\end{tabularx}
\end{table}

% ---------------- SÍNTESIS GENERAL ----------------
\section*{Síntesis general de la clase}
\addcontentsline{toc}{section}{Síntesis general de la clase}

La clase recorre un mismo problema desde ángulos sucesivos: cuánta libertad tiene una persona para organizar su actividad económica y cuánta responsabilidad le exige esa libertad a cambio. El punto de partida es un contexto de transformación acelerada del mundo del trabajo, que impulsa una reforma proyectada de la Ley General de Sociedades y que ya viene siendo anticipada por el Código Civil y Comercial a través del principio de autonomía de la voluntad. Sobre esa base, la clase repasa el sistema clásico de tipicidad societaria —con la Sociedad Anónima y la Sociedad de Responsabilidad Limitada como tipos de mayor uso— y muestra sus límites prácticos cuando se aplica sin criterio a estructuras pequeñas. Frente a esa rigidez, aparecen dos vías de flexibilización: la sociedad simple, que prescinde de un tipo previsto a cambio de un régimen de responsabilidad distinto, y la Sociedad Anónima Unipersonal, que —pese a haber sido pensada para el pequeño emprendedor— terminó exigiendo un capital mínimo elevado y sindicatura obligatoria, quedando reservada en la práctica a grandes sociedades vehículo. La Sociedad por Acciones Simplificada, creada por la Ley de Apoyo al Capital Emprendedor, ofrece hoy la alternativa más flexible y digital para el emprendedor real. El cierre conceptual de la clase —``la forma no hace el fondo''— conecta todos estos temas con una reflexión más amplia sobre el futuro de las profesiones jurídicas: lo que importa no es el soporte de un acto, sino la función que ese acto efectivamente cumple, tanto en el derecho societario como en el ejercicio profesional de quienes lo aplican.

\end{document}
